% =============================================================
%  RESUME TEMPLATE — moderncv banking style
%  -----------------------------------------------------------
%  How to use:
%    1. Replace every [YOUR_xxx] placeholder with your own data.
%    2. Section structure is ordered by methodology priority:
%       Summary  →  Experience (or Education-first for new grads)
%                →  Projects (optional)  →  Skills  →  Education
%    3. Keep ONE resume per role direction.
%       File-name your output: LastName_Resume_<Role>.pdf
%       e.g. Smith_Resume_AIEngineer.pdf
%    4. Build with:  pdflatex resume.tex
% =============================================================

\documentclass[11pt,a4paper,sans]{moderncv}

% Theme setup -------------------------------------------------
\moderncvstyle{banking}
\moderncvcolor{black}

% Encoding & layout packages ---------------------------------
\usepackage[utf8]{inputenc}
\usepackage[T1]{fontenc}
\usepackage{cmbright}
\renewcommand{\familydefault}{\sfdefault}
\usepackage{enumitem}
\usepackage{geometry}
\usepackage{ragged2e}
\usepackage[hidelinks]{hyperref}
\geometry{top=0.8cm, bottom=0.8cm, left=0.8cm, right=0.8cm}

% Custom helper for project entries (used in Projects section)
\newcommand{\projectentry}[2]{
 \vspace{-0.3em}
 \textbf{#1} \hfill #2\\[-1.0em]
}

% Layout tweaks ----------------------------------------------
\setlength{\hintscolumnwidth}{0pt}
\setlength{\maincolumnwidth}{\textwidth}
\setlength{\parindent}{0pt}

% =============================================================
%  PERSONAL INFO
%  -----------------------------------------------------------
%  \name takes {FirstName}{LastName}.
%  Drop any line you don't need (e.g. comment out \homepage if
%  you don't have a portfolio).
% =============================================================
\name{[YOUR_FIRST_NAME]}{[YOUR_LAST_NAME]}
\phone[mobile]{[YOUR_PHONE]}
\email{[YOUR_EMAIL]}
\social[linkedin]{[YOUR_LINKEDIN_HANDLE]}        % e.g. linkedin.com/in/yourname
\homepage{[YOUR_PORTFOLIO_OR_GITHUB]}             % optional — comment out if none
\extrainfo{[YOUR_CITY_STATE]}                     % e.g. Seattle, WA

% Itemize styling --------------------------------------------
\setlist[itemize]{
 leftmargin=*,
 label=\textbullet,
 labelsep=0.5em,
 itemsep=0.2em,
 partopsep=0pt,
 topsep=0.25em,
 parsep=0.25ex
}

\begin{document}
\makecvtitle
\vspace*{-2em}
\pagestyle{empty}

% =============================================================
%  SECTION 1 — SUMMARY  (a.k.a. "What I Do" / "Profile")
%  -----------------------------------------------------------
%  1–2 lines, ~25–40 words. Maximum information density.
%  Anchor to (a) what you do, (b) at what scale / with what proof,
%  (c) why it matters for THIS target role.
%
%  Tip: tailor this paragraph for every role direction.
%  This is the single most important section after recruiters
%  scan the top 1/3 of page 1.
% =============================================================
\vspace{-0.5em}
\section{Summary}
[YOUR_SUMMARY_PARAGRAPH — 1 to 2 lines, role-targeted, tied to a concrete proof point such as a deployed product URL, a measurable outcome, or a senior endorsement.]
\vspace{-0.6em}

% =============================================================
%  SECTION 2 — EXPERIENCE
%  -----------------------------------------------------------
%  Pack this section. Recruiters value Experience density above
%  Projects. Anything that has an advisor / team / company /
%  client behind it belongs here, not in Projects.
%
%  Section ordering rules:
%  - New grad / 0–1 YOE  → Education comes BEFORE Experience.
%  - 2+ YOE              → Experience comes FIRST (as below).
%
%  Optional: split into TWO sections when role-relevant work
%  would otherwise be buried by chronology, e.g.:
%    \section{AI Engineering Experience}   % role-direction first
%    \section{Professional Experience}     % everything else
%
%  Bullet-writing checklist:
%  - Lead with a strong verb (Architected / Shipped / Drove /
%    Led / Spearheaded / Productionized / Benchmarked / Refactored)
%  - Avoid weak verbs (Worked / Helped / Used / Participated /
%    Assisted)
%  - Limit total intensifier adverbs (Independently / Proactively /
%    Successfully / End-to-end) to 5 across the WHOLE resume
%  - Attach each bullet to at least one impact dimension:
%      Endorsement (who validated / who participated)
%      Scope       (team size, duration, budget, cross-team)
%      Marketing   (users, revenue, cost, latency, conversion)
% =============================================================

\section{Experience}

\textbf{[ROLE_TITLE], [COMPANY], [LOCATION]} \hfill \textbf{[START_DATE] -- [END_DATE]} \\
\vspace{-1.2em}
\begin{itemize}
  \item [STRONG_VERB] [WHAT_YOU_BUILT_OR_OWNED] on [TECH_STACK_OR_CONTEXT], [CONCRETE_DETAIL_THAT_PROVES_DEPTH], [QUANTIFIED_OUTCOME — Endorsement / Scope / Marketing].
  \item [STRONG_VERB] [SECOND_NOTABLE_DELIVERABLE], [HOW_YOU_DID_IT — name the technique / design choice], [QUANTIFIED_OUTCOME].
  \item [STRONG_VERB] [THIRD_DELIVERABLE — ideally a different impact dimension from the first two], [SPECIFIC_DETAIL], [QUANTIFIED_OUTCOME].
  \item [STRONG_VERB] [SCOPE_OR_LEADERSHIP_CONTRIBUTION — mentoring, cross-team alignment, hiring, etc.], [SUPPORTING_DETAIL].
\end{itemize}
\vspace{0.5em}

% --- Repeat the role block above for each prior role ---------

\textbf{[ROLE_TITLE], [COMPANY], [LOCATION]} \hfill \textbf{[START_DATE] -- [END_DATE]} \\
\vspace{-1.2em}
\begin{itemize}
  \item [BULLET_1]
  \item [BULLET_2]
  \item [BULLET_3]
\end{itemize}
\vspace{0.5em}

% Timeline overlap handling:
% If you held two roles concurrently, do NOT hide one by demoting
% it to Projects. Use explicit labels in the title line, e.g.:
%   "Visiting Researcher (Remote, Part-time)"
%   "Consulting Engagement (Part-time)"
%   "Concurrent Role"
% Make timing self-evident — don't make the hiring manager guess.

% =============================================================
%  SECTION 3 — PROJECTS  (OPTIONAL)
%  -----------------------------------------------------------
%  Keep to 2–3 max, only role-relevant ones.
%  Skip this section ENTIRELY for PM / Marketing / DA / DS-analytics
%  roles — push everything into Experience instead.
%  If a project has a deployed URL or an advisor / team behind it,
%  consider whether it belongs under Experience first.
%
%  Uncomment the block below to use it.
% =============================================================

% \section{Projects}
%
% \projectentry{[PROJECT_NAME] — [ONE_LINE_TAGLINE]}{[DATE_RANGE_OR_DEPLOYED_URL]}
% \begin{itemize}
%   \item [STRONG_VERB] [WHAT_YOU_BUILT], [TECH_AND_DESIGN_DETAILS], [OUTCOME_OR_VALIDATION].
%   \item [STRONG_VERB] [TECHNICAL_DEPTH_DETAIL — show production-grade thinking, not toy demo].
%   \item [STRONG_VERB] [DIFFERENTIATING_DETAIL — eval harness, perf benchmark, scaling result, etc.].
% \end{itemize}
% \vspace{0.5em}

% =============================================================
%  SECTION 4 — TECHNICAL SKILLS
%  -----------------------------------------------------------
%  Reorganize the categories per target role. The category that
%  the JD weights most heavily must come FIRST.
%
%  Examples of role-direction-driven ordering:
%
%    AI Engineer:
%      Languages → ML/AI Systems → AI Application → Backend → Cloud → Frontend
%
%    Data Scientist:
%      Languages → Statistics/ML → Experimentation/Causal → Data Tools → Cloud
%
%    Product Manager:
%      (often skip this section; if kept: Tools → Methods → Data Stack)
%
%    UX Designer:
%      Tools (Figma, etc.) → Methods → Research → Frontend basics
%
%  Keyword coverage target: ≥ 75% of high-frequency terms from
%  the 2–3 representative JDs you sampled for this role.
% =============================================================
\section{Technical Skills}
\vspace{-0.2em}
\noindent\textbf{[CATEGORY_1 — most JD-relevant]}: [comma-separated skill list]\par
\noindent\textbf{[CATEGORY_2]}: [comma-separated skill list]\par
\noindent\textbf{[CATEGORY_3]}: [comma-separated skill list]\par
\noindent\textbf{[CATEGORY_4]}: [comma-separated skill list]\par
\noindent\textbf{[CATEGORY_5]}: [comma-separated skill list]\par
\vspace{-0.2em}

% =============================================================
%  SECTION 5 — EDUCATION
%  -----------------------------------------------------------
%  New grads / 0–1 YOE: place this section BEFORE Experience.
%  2+ YOE: keep it here, near the bottom.
%
%  Skip Coursework unless you're a new grad with directly
%  relevant courses. GPA: include only if ≥ 3.5/4.0 or ≥ 80/100.
%
%  Add scholarships, awards, or thesis title only if directly
%  relevant to the target role.
% =============================================================
\section{Education}
\vspace{-0.2em}
\textbf{[UNIVERSITY_1]} \hfill {[DATE_RANGE_1]} \\
\textit{[DEGREE_1 — e.g. M.S. in Computer Science, Specialization: ...]} \\[0.5em]

\textbf{[UNIVERSITY_2]} \hfill {[DATE_RANGE_2]} \\
\textit{[DEGREE_2]} \\[0.5em]

% \textbf{[UNIVERSITY_3]} \hfill {[DATE_RANGE_3]} \\
% \textit{[DEGREE_3]} \\
\vspace{-0.2em}

% =============================================================
%  OPTIONAL SECTIONS
%  -----------------------------------------------------------
%  Add only when role-relevant and the addition strengthens
%  candidacy. Examples:
%
%  \section{Selected Publications}
%  \section{Patents}
%  \section{Open Source}
%  \section{Speaking & Teaching}
%  \section{Certifications}            % AWS / GCP / PMP / CFA only when relevant
%  \section{Volunteer Leadership}      % only when it demonstrates Scope/Endorsement
%
%  Hard skip — do NOT include:
%  - Hobbies (unless a senior creative role asks for it)
%  - References (always "available upon request" implicit)
%  - Photo (illegal/optional in US; risks bias-screening)
%  - Date of birth, marital status, nationality (US convention)
% =============================================================

\end{document}
